%!TEX root = ../../thesis.tex

\section{Relation to Prior Work}\label{sec:corpus-prior-work}

% TODO: This formulation is a direct descendant of \cite{bottinger_deep_2018}, so give
% it a proper comparison table --- one row per MDP component, one column per approach.
% Rows: RL algorithm, state, action, mutation randomness, reward, corpus, horizon,
% target and instrumentation, reported result.
%
% What is inherited: the MDP formulation of fuzzing itself --- states as input
% fragments, actions as rewrite rules, reward from runtime coverage. This thesis
% extends the canonical formulation rather than proposing a new one; say that plainly.
%
% What differs, and why each difference is forced rather than cosmetic:
%   1. model-based planning instead of model-free DQN, which forces the determinism
%      of \cref{sec:corpus-determinism},
%   2. a learned corpus --- selection *and* retention/eviction --- instead of a single
%      seed; note that the multi-seed extension is mentioned there as future work with
%      a *random* choice of seed, never learned and never evaluated. This is the novel
%      core of the chapter,
%   3. a long horizon with seed lineage instead of a reset after every mutation.
%
% Also inherit the cautions, and state them before presenting any result:
%   - the reported gain over a random baseline was modest and called preliminary, so
%     calibrate expectations about beating random fuzzing,
%   - learned control of the mutation *offset* did not beat a random offset there,
%     which puts the learned position factor of \cref{sec:corpus-action-space} under
%     suspicion --- \cref{sec:results-ablations} should test whether it earns its
%     contribution to the branching factor.
%
% Finally, relate to the seed-scheduling line of \cref{sec:rlfuzz-prior-work}
% (bandit and MCTS-over-seed-tree approaches): those learn scheduling around a
% conventional fuzzer, whereas here selection, mutation and retention are one joint
% action in a single MDP.
